\origpage{27}{}
\textbf{6.6} \textit{For all integers $p \ge 0$ and $q \ge 1$ there exists $\lambda \ge 0$ such that, if $G$ has a presentation in a disc $\Sigma$ of depth $\le q$ and defect $3$, and there is an exhaustive vital $p$-linkage \fix{in} $G$, then there are at most $\lambda$ $\Sigma$-jumps.}\fixnote{fix:6.6}{The word ``in'' is missing, as the statements of (6.4) and (6.5) show.}

\textbf{Proof.} Choose $\lambda' \ge 0$ so that (6.5) holds with $p, q, \lambda$ replaced by $p + 2q, q, \lambda'$, and let $\lambda = 3\lambda'$. We claim that $\lambda$ satisfies (6.6). For let $\Gamma, F_1^*, F_2^*, F_3^*, H_1, H_2, \delm{H_2}\fix{H_3}, \mu$\fixnote{fix:6.6-proof-H3}{The third graph is $H_3$.} be a presentation of $G$ in the disc $\Sigma$, of depth $\le q$, and let there be an exhaustive vital $p$-linkage. For $i = 1, 2, 3$, choose $v_i \in V(\Gamma) \cap F_i^*$ such that there is a region $r$ of $\Gamma$ in $\Sigma$ incident with $v_1, v_2$ and $v_3$. (It is an easy exercise to prove that this choice is possible, since $V(\Gamma) \cap bd(\Sigma) \subseteq F_1^* \cup F_2^* \cup F_3^*$ \fix{and, because $L$ is exhaustive, $\Gamma$ is a forest in which every vertex not in $bd(\Sigma)$ has degree 2}.)\fixnote{fix:6.6-proof-region}{The reason given is not enough: if $\Gamma$ is a star with its center inside the disc and one leaf in each $F_i^*$, no region is incident with all three leaves. The choice is possible because $L$ is exhaustive: as stated on page 24, $\Gamma$ is then a forest in which every vertex not in $bd(\Sigma)$ has degree 2, which excludes the star. The same fact is used in Claim (3) in the proof of (6.3).} Choose $v_0 \in r \setminus bd(\Sigma)$, and for $i = 1, 2, 3$ let $F_i \subseteq r \cup \{v_i\}$ be a line with ends $v_0$ and $v_i$, so that $F_1, F_2$ and $F_3$ are mutually disjoint except for $v_0$, and $F_i \cap bd(\Sigma) = \{v_i\}$ $(i = 1, 2, 3)$. For $i = 1, 2, 3$ let $J_i \subseteq bd(\Sigma)$ be the support with ends $\{v_1, v_2, v_3\} \setminus \{v_i\}$ which does not include $F_i^*$. Let $\Delta_1 \subseteq \Sigma$ be the closed disc bounded by $J_1 \cup F_2 \cup F_3$ and define $\Delta_2, \Delta_3$ similarly. Now
\begin{equation*}
\Gamma \cap \Delta_1, F_2^* \cap J_1, F_3^* \cap J_1, \mu(F_2^* \cap J_1), \mu(F_3^* \cap J_1)
\end{equation*}
and the restriction of $\mu$ to $(F_2^* \cup F_3^*) \cap J_1$ is a presentation of $\Gamma \cap \Delta_1 \cup \mu(F_2^* \delm{\cup}\fix{\cap} J_1) \cup \mu(F_3^* \delm{\cup}\fix{\cap} J_1)$\fixnote{fix:6.6-proof-Delta1}{The graphs at $\Delta_1$ are $\mu(F_2^* \cap J_1)$ and $\mu(F_3^* \cap J_1)$, as in the display just before. The map $\mu$ is defined on the parts of the supports, and $F_2^* \cup J_1$ is not a part of $F_2^*$.} of depth $\le q$ and defect $2$, and this graph has an exhaustive vital $(p + 2q)$-linkage. Consequently, $\Delta_1$ includes at most $\lambda'$ $\Sigma$-jumps. Similarly so do $\Delta_2$ and $\Delta_3$, and since every $\Sigma$-jump belongs to one of these discs, it follows that there are at most $3\lambda' = \lambda$ $\Sigma$-jumps altogether, as required. $\square$

\textbf{6.7} \textit{For all integers $p, r \ge 0$ and $q \ge 1$ there exists $\lambda \ge 0$ such that, if $G$ has a presentation in a disc $\Sigma$ of depth $\le q$ and defect $\le r$, and there is an exhaustive vital $p$-linkage in $G$, then there are at most $\lambda$ $\Sigma$-jumps.}

\textbf{Proof.} We prove the result for all $p$ and $q$ by induction on $r$. By (6.4), (6.5), (6.6) we may assume that $r \ge 4$, and for all $p' \ge 0$ and $q' \ge 1$ there exists $\lambda(p', q') \ge 0$ such that, if $G$ has a presentation in a disc $\Sigma$ of depth $\le q'$ and defect $< r$, and there is an exhaustive vital $p'$-linkage in $G$, then there are at most $\lambda(p', q')$ $\Sigma$-jumps. Let $\lambda = 2\lambda(p + 2q, q)$. We claim that $\lambda$ satisfies (6.7). For let $\Gamma, F_1^*, \dots, F_r^*, H_1, \dots, H_r, \mu$ be a presentation of $G$ in the disc $\Sigma$ of depth $\le q$ and let there be an exhaustive vital $p$-linkage in $G$. Choose a region $r$ incident with vertices in at least three of $F_1^*, \dots, F_r^*$, as in the proof of (6.6). Since two of these three are non-consecutive, because $r \ge 4$, it follows that there is a proper line $F \subseteq \Sigma$ with ends $a, b \in V(\Gamma)$, and with $F \cap U(\Gamma) = \{a, b\}$, such that both of the lines $J_1, J_2 \subseteq bd(\Sigma)$ with ends $a$ and $b$ are disjoint from at least one of $F_1^*, \dots, F_r^*$.
